Galaxy deblending is necessary when overlapping astronomical sources bias
measurements of flux, color, morphology, or downstream redshift inference. Real
survey deblending is difficult because overlapping light can be physically
ambiguous and because imaging conditions introduce sky, PSF, detector, and
crowding effects.

This work focuses on a narrower controlled question: given synthetic blends
with known clean targets, can compact U-Net models recover target galaxies more
accurately than simple baselines, and how does performance change under harder
blend conditions? This controlled setting does not solve survey-grade
deblending, but it makes reconstruction targets known and enables direct
evaluation of model behavior.

The experiments compare a progression of learned formulations. Thayer-Direct
predicts the clean target from the blended image. Thayer-Residual predicts the
contaminant residual layer and reconstructs by subtraction. Thayer-BR v0.1
keeps the residual objective but changes the training distribution toward
normal blends, core-overlap blends, and brightness/size stress blends.
Thayer-BR v0.2 Moderate keeps the same residual family and adds a moderate
affected/core-weighted residual loss. Thayer-BR v0.2 Strong is evaluated as a
loss-weighting ablation.
